% Reviewed translation: PDF pages 76--81 / printed pages 62--67.

\MathBookSourcePage{76}
通常把 $J$ 称为与问题 $(P)$ 相联系的能量泛函,
把 $\mathcal L$ 称为与问题 $(Q)$ 相联系的 Lagrange 泛函.
考虑如下问题, 称为问题 $(L)$: 求 Lagrange 泛函 $\mathcal L$ 在
$X\times M$ 上的鞍点 $(u,\lambda)\in X\times M$, 即求
$(u,\lambda)\in X\times M$, 使
\begin{equation}\label{eq:abstract-saddle-point}
\mathcal L(u,\mu)\le\mathcal L(u,\lambda)\le\mathcal L(v,\lambda)
\qquad\forall v\in X,\quad\forall\mu\in M.
\end{equation}

\setcounter{statement}{1}
\begin{Theorem}\label{thm:abstract-saddle-equivalence}
假设 Theorem~\ref{thm:abstract-well-posedness} 的条件 (i) 和 (ii) 成立.
再假设双线性形式 $a(\cdot,\cdot)$ 在 $X$ 上对称且半正定, 即
\begin{equation}\label{eq:abstract-semipositive}
a(v,v)\ge0
\qquad\forall v\in X.
\end{equation}
则问题 $(L)$ 有唯一解 $(u,\lambda)\in X\times M$,
它恰好是问题 $(Q)$ 的解.
\end{Theorem}
\begin{Proof}
\eqref{eq:abstract-saddle-point} 中第一个不等式可写成
\[
b(u,\mu-\lambda)\le\langle\chi,\mu-\lambda\rangle
\qquad\forall\mu\in M.
\]
由于 $\mu$ 是 $M$ 的任意元素, 这等价于
\[
b(u,\mu)=\langle\chi,\mu\rangle
\qquad\forall\mu\in M.
\]
其次, \eqref{eq:abstract-saddle-point} 中第二个不等式等价于
\[
\mathcal L(u,\lambda)=\inf_{v\in X}\mathcal L(v,\lambda).
\]
根据假设, $a(\cdot,\cdot)$ 对称, 所以
\[
D_v\mathcal L(u,\lambda)\mathbin{:}v
=a(u,v)+b(v,\lambda)-\langle l,v\rangle.
\]
此外, 利用 
\eqref{eq:abstract-semipositive}, 有
\[
D_{vv}^2\mathcal L(u,\lambda)\mathbin{:}(v,v)
=a(v,v)\ge0
\qquad\forall v\in X.
\]
因此 $v\mapsto\mathcal L(v,\lambda)$ 是凸泛函, 其极小点 $u$ 由条件
$D_v\mathcal L(u,\lambda)\mathbin{:}v=0$ 刻画, 即
\[
a(u,v)+b(v,\lambda)=\langle l,v\rangle
\qquad\forall v\in X.
\]
所以 $(u,\lambda)$ 是 $\mathcal L$ 的鞍点, 当且仅当它也是问题 $(Q)$ 的解.
Theorem 得证.
\end{Proof}

\setcounter{statement}{1}
\begin{Corollary}\label{cor:abstract-saddle-minmax}
在 Theorem~\ref{thm:abstract-saddle-equivalence} 的假设下,
问题 $(Q)$ 的解 $(u,\lambda)$ 由下式刻画:
\begin{equation}\label{eq:abstract-saddle-minmax}
\operatorname{Min}_{v\in X}\left(\sup_{\mu\in M}\mathcal L(v,\mu)\right)
=\mathcal L(u,\lambda)
=\operatorname{Max}_{\mu\in M}\left(\inf_{v\in X}\mathcal L(v,\mu)\right),
\end{equation}
其中用 $\operatorname{Min}$ 代替 $\inf$(相应地, 用 $\operatorname{Max}$ 代替 $\sup$)表示极值可以达到.
\end{Corollary}
\begin{Proof}
这是一般优化结果的推论: Lagrange 泛函 $\mathcal L$ 有鞍点
$(u,\lambda)$, 当且仅当
\begin{equation}\tag{\ref*{eq:abstract-saddle-minmax}$'$}
\label{eq:abstract-saddle-minmax-prime}
\operatorname{Min}_{v\in X}\left(\sup_{\mu\in M}\mathcal L(v,\mu)\right)
=\operatorname{Max}_{\mu\in M}\left(\inf_{v\in X}\mathcal L(v,\mu)\right),
\end{equation}
并且这个量等于 $\mathcal L(u,\lambda)$.

为方便读者, 回顾其证明. 置
\[
\phi(v)=\sup_{\mu\in M}\mathcal L(v,\mu),
\qquad
\phi^*(\mu)=\inf_{v\in X}\mathcal L(v,\mu).
\]
先注意
\begin{equation}\label{eq:abstract-weak-duality}
\sup_{\mu\in M}\phi^*(\mu)
\le\inf_{v\in X}\phi(v).
\end{equation}
事实上, 因为
\[
\mathcal L(v,\mu)\le\phi(v)
\qquad\forall v\in X,\quad\forall\mu\in M,
\]
所以
\[
\phi^*(\mu)\le\inf_{v\in X}\phi(v)
\qquad\forall\mu\in M,
\]
从而得到 
\eqref{eq:abstract-weak-duality}.

\MathBookSourcePage{77}
现在设 $(u,\lambda)$ 是 $\mathcal L$ 的鞍点. 这意味着
\[
\phi(u)=\mathcal L(u,\lambda)=\phi^*(\lambda).
\]
因此
\begin{equation}\label{eq:abstract-saddle-dual-chain}
\inf_{v\in X}\phi(v)
\le\phi(u)=\phi^*(\lambda)
\le\sup_{\mu\in M}\phi^*(\mu).
\end{equation}
结合 
\eqref{eq:abstract-weak-duality}, 得
\[
\inf_{v\in X}\phi(v)=\phi(u)=\mathcal L(u,\lambda)
=\phi^*(\lambda)=\sup_{\mu\in M}\phi^*(\mu),
\]
特别地, 这蕴含 
\eqref{eq:abstract-saddle-minmax-prime}.

反之, 假设 
\eqref{eq:abstract-saddle-minmax-prime} 成立, 并设 $u$ 达到 $\phi$ 的极小值,
$\lambda$ 达到 $\phi^*$ 的极大值. 此时
\eqref{eq:abstract-saddle-minmax-prime} 表明
\[
\phi(u)=\phi^*(\lambda).
\]
另一方面, $\phi$ 和 $\phi^*$ 的定义蕴含
\[
\phi^*(\lambda)\le\mathcal L(u,\lambda)\le\phi(u).
\]
所以
\[
\phi(u)=\mathcal L(u,\lambda)=\phi^*(\lambda),
\]
即 $(u,\lambda)$ 是 $\mathcal L$ 的鞍点.
最后, 这个优化结果与问题 $(Q)$ 解的唯一性给出刻画
\eqref{eq:abstract-saddle-minmax}.
\end{Proof}

\MathBookSourcePage{78}
注意
\[
\sup_{\mu\in M}\mathcal L(v,\mu)
=\begin{cases}
+\infty,&v\notin V(\chi),\\
J(v),&v\in V(\chi).
\end{cases}
\]
因此
\[
\inf_{v\in X}\sup_{\mu\in M}\mathcal L(v,\mu)
=\inf_{v\in V(\chi)}J(v).
\]
因为 $b(u,\mu)=\langle\chi,\mu\rangle$,
\eqref{eq:abstract-saddle-minmax} 中第一个等式变为
\begin{equation}\label{eq:abstract-constrained-minimization}
J(u)=\operatorname{Min}_{v\in V(\chi)}J(v).
\end{equation}
所以问题 $(P)$ 的解 $u$ 可刻画为约束优化问题
\eqref{eq:abstract-constrained-minimization} 的解.

另一方面, 考虑 
\eqref{eq:abstract-constrained-minimization} 的对偶问题
\[
\operatorname{Max}_{\mu\in M}\inf_{v\in X}\mathcal L(v,\mu).
\]
除 Theorem~\ref{thm:abstract-saddle-equivalence} 的假设外,
再假设双线性形式 $a(\cdot,\cdot)$ 是 $X$-椭圆的, 即
\begin{equation}\label{eq:abstract-x-ellipticity}
a(v,v)\ge\alpha\|v\|_X^2
\qquad\forall v\in X,\quad\alpha>0.
\end{equation}
对任意 $\mu\in M$, 定义 $u(\mu)\in X$ 为如下优化问题的解:
\[
\mathcal L(u(\mu),\mu)=\inf_{v\in X}\mathcal L(v,\mu).
\]
在 Theorem~\ref{thm:abstract-saddle-equivalence} 的证明中已经看到,
这个问题等价于方程
\begin{equation}\label{eq:abstract-dual-state-equation}
a(u(\mu),v)=\langle l,v\rangle-b(v,\mu)
\qquad\forall v\in X.
\end{equation}
由 
\eqref{eq:abstract-x-ellipticity}, 问题
\eqref{eq:abstract-dual-state-equation} 确有唯一解 $u(\mu)\in X$.
由于 $u(\lambda)=u$, 注意到
\[
\mathcal L(u(\lambda),\lambda)
=\operatorname{Max}_{\mu\in M}\mathcal L(u(\mu),\mu).
\]

给出 $\mathcal L(u(\mu),\mu)$ 的简单表达式. 有
\[
\mathcal L(u(\mu),\mu)
=(1/2)a(u(\mu),u(\mu))-\langle l,u(\mu)\rangle
+b(u(\mu),\mu)-\langle\chi,\mu\rangle,
\]
所以由 
\eqref{eq:abstract-dual-state-equation},
\[
\mathcal L(u(\mu),\mu)
=-(1/2)a(u(\mu),u(\mu))-\langle\chi,\mu\rangle.
\]
于是置
\begin{equation}\label{eq:abstract-dual-functional-k}
K(\mu)=(1/2)a(u(\mu),u(\mu))+\langle\chi,\mu\rangle,
\end{equation}
则 
\eqref{eq:abstract-constrained-minimization} 的对偶问题变为

\MathBookSourcePage{79}
\begin{equation}\label{eq:abstract-dual-minimization}
K(\lambda)=\operatorname{Min}_{\mu\in M}K(\mu).
\end{equation}
因此问题 $(Q)$ 的解的第二个分量 $\lambda$ 可刻画为无约束优化问题
\eqref{eq:abstract-dual-minimization} 的解.

\setcounter{subsection}{2}
\subsection{正则化或罚方法逼近}\label{subsec:abstract-regularization-penalty}

我们希望引入问题 $(Q)$ 的一个扰动形式, 它在实际中可能更容易求解.
除 $a(\cdot,\cdot)$ 和 $b(\cdot,\cdot)$ 外, 再考虑第三个连续双线性形式
\[
c(\cdot,\cdot):M\times M\longrightarrow\mathbb R,
\]
其范数为
\[
\|c\|=\sup_{\substack{\lambda,\mu\in M\\\lambda,\mu\ne0}}
\frac{c(\lambda,\mu)}{\|\lambda\|_M\|\mu\|_M}.
\]
假设形式 $c(\cdot,\cdot)$ 是 $M$-椭圆的, 即存在常数 $\gamma>0$, 使
\begin{equation}\label{eq:abstract-c-ellipticity}
c(\mu,\mu)\ge\gamma\|\mu\|_M^2
\qquad\forall\mu\in M.
\end{equation}
对形式 $c(\cdot,\cdot)$, 定义算子 $C\in\mathcal L(M;M')$:
\begin{equation}\label{eq:abstract-operator-c}
\langle C\lambda,\mu\rangle=c(\lambda,\mu)
\qquad\forall\lambda,\mu\in M.
\end{equation}
设 $\varepsilon>0$ 是趋于零的参数. 考虑如下问题, 称为问题 $(Q^\varepsilon)$:
求 $(u^\varepsilon,\lambda^\varepsilon)\in X\times M$, 使
\begin{equation}\label{eq:abstract-regularized-q-first}
a(u^\varepsilon,v)+b(v,\lambda^\varepsilon)=\langle l,v\rangle
\qquad\forall v\in X,
\end{equation}
\begin{equation}\label{eq:abstract-regularized-q-second}
-\varepsilon c(\lambda^\varepsilon,\mu)+b(u^\varepsilon,\mu)
=\langle\chi,\mu\rangle
\qquad\forall\mu\in M.
\end{equation}
因为算子 $C$ 非奇异, 方程 
\eqref{eq:abstract-regularized-q-second} 等价于
\begin{equation}\label{eq:abstract-lambda-epsilon}
\lambda^\varepsilon=(1/\varepsilon)C^{-1}(Bu^\varepsilon-\chi).
\end{equation}
在 
\eqref{eq:abstract-regularized-q-first} 中代入 $\lambda^\varepsilon$,
得到如下问题, 称为问题 $(P^\varepsilon)$, 它显然与问题 $(Q^\varepsilon)$ 等价:
求 $u^\varepsilon\in X$, 使
\begin{equation}\label{eq:abstract-penalty-problem}
a(u^\varepsilon,v)+(1/\varepsilon)
\langle C^{-1}Bu^\varepsilon,Bv\rangle
=\langle l,v\rangle+(1/\varepsilon)\langle C^{-1}\chi,Bv\rangle
\qquad\forall v\in X.
\end{equation}

\setcounter{statement}{2}
\begin{Theorem}\label{thm:abstract-penalty-error}
假设条件 
\eqref{eq:abstract-inf-sup} 和
\eqref{eq:abstract-c-ellipticity} 成立. 再假设存在常数 $\alpha>0$, 使
\begin{equation}\label{eq:abstract-combined-ellipticity}
a(v,v)+\langle C^{-1}Bv,Bv\rangle
\ge\alpha\|v\|_X^2
\qquad\forall v\in X.
\end{equation}
则当 $\varepsilon\le1$ 时, 问题 $(Q)$ 与 $(Q^\varepsilon)$ 分别在
$X\times M$ 中有唯一解 $(u,\lambda)$ 和
$(u^\varepsilon,\lambda^\varepsilon)$. 此外, 当
$\varepsilon\le\varepsilon_0$ 足够小时, 有误差估计
\begin{equation}\label{eq:abstract-penalty-error-bound}
\|u^\varepsilon-u\|_X+\|\lambda^\varepsilon-\lambda\|_M
\le K\varepsilon(\|l\|_{X'}+\|\chi\|_{M'}),
\end{equation}
其中常数 $K$ 只依赖于 $\alpha$, $\beta$, $\|a\|$, $\|b\|$ 和 $\|c\|$.
\end{Theorem}

\MathBookSourcePage{80}
\begin{Proof}
假设 
\eqref{eq:abstract-combined-ellipticity} 蕴含双线性形式
$a(\cdot,\cdot)$ 是 $V$-椭圆的. 因而由
Corollary~\ref{cor:abstract-v-elliptic-well-posedness},
问题 $(Q)$ 有唯一解 $(u,\lambda)\in X\times M$.

由 
\eqref{eq:abstract-c-ellipticity} 和
\eqref{eq:abstract-combined-ellipticity}, 双线性形式
\[
(u,v)\longmapsto a(u,v)+(1/\varepsilon)
\langle C^{-1}Bu,Bv\rangle
\]
在 $\varepsilon\le1$ 时是 $X$-椭圆的. 所以问题 $(P^\varepsilon)$
在 $X$ 中恰有一个解 $u^\varepsilon$. 因而, 若用
\eqref{eq:abstract-lambda-epsilon} 定义 $\lambda^\varepsilon$,
则 $(u^\varepsilon,\lambda^\varepsilon)$ 是问题 $(Q^\varepsilon)$ 的唯一解.

还需建立估计 
\eqref{eq:abstract-penalty-error-bound}. 由
\eqref{eq:abstract-q-first} 与
\eqref{eq:abstract-regularized-q-first}, 以及
\eqref{eq:abstract-q-second} 与
\eqref{eq:abstract-regularized-q-second}, 得
\begin{equation}\label{eq:abstract-penalty-error-system}
\left\{
\begin{aligned}
a(u^\varepsilon-u,v)+b(v,\lambda^\varepsilon-\lambda)&=0
&&\forall v\in X,\\
-\varepsilon c(\lambda^\varepsilon,\mu)
+b(u^\varepsilon-u,\mu)&=0
&&\forall\mu\in M.
\end{aligned}
\right.
\end{equation}
第一个方程 
\eqref{eq:abstract-penalty-error-system} 与
\eqref{eq:abstract-inf-sup} 给出
\[
\beta\|\lambda^\varepsilon-\lambda\|_M
\le\sup_{v\in X}\frac{b(v,\lambda^\varepsilon-\lambda)}{\|v\|_X}
\le\|a\|\|u^\varepsilon-u\|_X,
\]
从而
\begin{equation}\label{eq:abstract-penalty-lambda-bound}
\|\lambda^\varepsilon-\lambda\|_M
\le(\|a\|/\beta)\|u^\varepsilon-u\|_X.
\end{equation}
接着在 
\eqref{eq:abstract-penalty-error-system} 中取
$v=u^\varepsilon-u$ 和 $\mu=\lambda^\varepsilon-\lambda$, 得
\[
\begin{aligned}
a(u^\varepsilon-u,u^\varepsilon-u)
&=-\varepsilon c(\lambda,\lambda^\varepsilon-\lambda)
-\varepsilon c(\lambda^\varepsilon-\lambda,
\lambda^\varepsilon-\lambda)\\
&\le-\varepsilon c(\lambda,\lambda^\varepsilon-\lambda).
\end{aligned}
\]
于是 
\eqref{eq:abstract-penalty-lambda-bound} 给出
\[
a(u^\varepsilon-u,u^\varepsilon-u)
\le\varepsilon(\|a\|\|c\|/\beta)
\|\lambda\|_M\|u^\varepsilon-u\|_X.
\]
另一方面, 由 
\eqref{eq:abstract-penalty-error-system} 中第二个方程,
\[
B(u^\varepsilon-u)=\varepsilon C\lambda^\varepsilon.
\]
因此
\[
\begin{aligned}
\langle C^{-1}B(u^\varepsilon-u),B(u^\varepsilon-u)\rangle
&=\varepsilon^2\langle C\lambda^\varepsilon,
\lambda^\varepsilon\rangle\\
&=\varepsilon^2c(\lambda^\varepsilon,\lambda^\varepsilon)
\le\varepsilon^2\|c\|\|\lambda^\varepsilon\|_M^2,
\end{aligned}
\]
并由 
\eqref{eq:abstract-penalty-lambda-bound},
\[
\langle C^{-1}B(u^\varepsilon-u),B(u^\varepsilon-u)\rangle
\le\varepsilon^2\|c\|
\{(\|a\|/\beta)\|u^\varepsilon-u\|_X+\|\lambda\|_M\}^2.
\]
所以假设 
\eqref{eq:abstract-combined-ellipticity} 给出形如
\[
\alpha x^2\le\varepsilon^2C_1(\|\lambda\|_M+x)^2
+\varepsilon C_2\|\lambda\|_Mx,
\]
的不等式, 其中 $x=\|u^\varepsilon-u\|_X$.

\MathBookSourcePage{81}
当 $\varepsilon$ 足够小时, 这等价于
\[
\|u^\varepsilon-u\|_X\le\varepsilon C_3\|\lambda\|_M.
\]
结合 
\eqref{eq:abstract-penalty-lambda-bound}, 即证明
\eqref{eq:abstract-penalty-error-bound}.
\end{Proof}

\setcounter{statement}{2}
\begin{Remark}\label{rem:abstract-regularization-without-infsup}
即使 inf-sup 条件 
\eqref{eq:abstract-inf-sup} 不成立, 问题 $(Q^\varepsilon)$ 仍有唯一解.
因此在 
\eqref{eq:abstract-regularized-q-second} 中,
$-\varepsilon c(\lambda^\varepsilon,\mu)$ 起正则化项的作用:
问题 $(Q^\varepsilon)$ 可视为问题 $(Q)$ 的正则化版本.
\end{Remark}

\setcounter{statement}{3}
\begin{Remark}\label{rem:abstract-penalized-functional}
当双线性形式 $a(\cdot,\cdot)$ 和 $c(\cdot,\cdot)$ 对称时,
双线性形式
\[
(u,v)\longmapsto a(u,v)+(1/\varepsilon)
\langle C^{-1}Bu,Bv\rangle
\]
也对称. 此时置
\[
J_\varepsilon(v)=J(v)+\frac1{2\varepsilon}
\langle C^{-1}(Bv-\chi),Bv-\chi\rangle,
\]
并利用 $\varepsilon\le1$ 时的椭圆性
\eqref{eq:abstract-combined-ellipticity}, 得到问题 $(P^\varepsilon)$
等价于求 $u^\varepsilon\in X$, 使
\[
J_\varepsilon(u^\varepsilon)=\inf_{v\in X}J_\varepsilon(v).
\]
这样, 我们用罚泛函 $J_\varepsilon$ 的无约束优化问题逼近了能量泛函 $J$
的约束优化问题 
\eqref{eq:abstract-constrained-minimization}, 其中
$\langle C^{-1}(Bv-\chi),Bv-\chi\rangle/(2\varepsilon)$
是与约束 $Bv=\chi$ 对应的罚项. 因而问题 $(P^\varepsilon)$ 是问题 $(P)$ 的罚方法版本.
\end{Remark}

事实上, 可以得到比 Theorem~\ref{thm:abstract-penalty-error} 更精确的结果,
即 $(u^\varepsilon,\lambda^\varepsilon)$ 关于 $\varepsilon$ 的幂的渐近展开.
按如下方式归纳定义序列 $(u_n,\lambda_n)\in X\times M$.
置 $\lambda_0=\lambda$, 并对任意整数 $n\ge1$, 令
$(u_n,\lambda_n)\in X\times M$ 为下列问题的解:
\begin{equation}\label{eq:abstract-asymptotic-recursion}
\left\{
\begin{aligned}
a(u_n,v)+b(v,\lambda_n)&=0
&&\forall v\in X,\\
b(u_n,\mu)&=c(\lambda_{n-1},\mu)
&&\forall\mu\in M.
\end{aligned}
\right.
\end{equation}
已知 $\lambda_{n-1}$ 后, 由
\eqref{eq:abstract-inf-sup},
\eqref{eq:abstract-combined-ellipticity} 和
Corollary~\ref{cor:abstract-v-elliptic-well-posedness}, 方程
\eqref{eq:abstract-asymptotic-recursion} 确实定义唯一的
$(u_n,\lambda_n)\in X\times M$.

\setcounter{statement}{3}
\begin{Theorem}\label{thm:abstract-penalty-asymptotic-expansion}
假设 Theorem~\ref{thm:abstract-penalty-error} 的条件成立.
则对所有整数 $N\ge1$, 当 $\varepsilon\le\varepsilon_0$ 足够小时,
\begin{equation}\label{eq:abstract-penalty-asymptotic-bound}
\left\|u^\varepsilon-u-\sum_{n=1}^N\varepsilon^nu_n\right\|_X
+\left\|\lambda^\varepsilon-\lambda-
\sum_{n=1}^N\varepsilon^n\lambda_n\right\|_M
\le K_N\varepsilon^{N+1}(\|l\|_{X'}+\|\chi\|_{M'}),
\end{equation}
其中常数 $K_N$ 只依赖于 $N$, $\alpha$, $\beta$, $\|a\|$, $\|b\|$ 和 $\|c\|$.
\end{Theorem}
